\section{The Cusp}

The flat three-dimensional boundary the bulk grows toward is the cusp, and the cusp is the ordinary three-dimensional space we live in. Physical space is not the curved bulk but its flat ideal boundary, one dimension down, and the fourth bulk dimension is the radial depth beneath each point of the cusp.

The cusp is a horosphere, a flat slice of the hyperbolic interior, and it is the cubic honeycomb, plain flat three-dimensional space. Its own growth is polynomial, like an ordinary Euclidean region, in sharp contrast to the exponential swelling of the bulk. This is why the dimension of physical space can be read directly off the cusp: the cusp grows like radius cubed, an exact three, and the bulk is one higher, a four. The same reading on the lower sibling gives a two-dimensional cusp and a three-dimensional bulk, the dimension-down control that fixes the meaning of the four.

Two things about physical space follow from the cusp being a flat aperiodic slice rather than a regular grid. First, it restores isotropy with scale. The slice has no repeating grain, so its roughness is random and averages away the farther out you look, while a flat repeating lattice keeps its facets at every scale. So a curved aperiodic substrate looks smooth and directionless at large scale, exactly what relativity needs, and a flat mesh cannot. Second, gravity on the cusp falls off the dimension-correct way. In three flat dimensions a static influence spreads over a surface that grows as the distance squared, so the Newtonian potential falls as one over the distance, while in two dimensions it would be a logarithm. The cusp being three-dimensional and flat forces the one-over-r law of gravity we measure.
